\documentclass[11pt]{article}
\usepackage[a4paper,margin=0.9in]{geometry}
\usepackage{amsmath,amssymb,amsthm,hyperref}
\newtheorem{theorem}{Theorem}
\newtheorem{lemma}{Lemma}
\title{An alternating Tur\'an inequality for normalized Euler numbers}
\author{Research proof and Lean formalization}
\date{October 27, 2026}
\begin{document}
\maketitle
\begin{abstract}
We give a proof of Sokal's Conjecture 3.3 for every starting index and every pair of positive shifts. Classical partial fractions express the normalized Euler numbers as a geometric factor times an alternating perturbation of one. A uniform contraction of the positive perturbation under a shift of two, followed by four finite sign cases, proves the conjecture. The actual generating function and all coefficient bridges are formalized in Lean.
\end{abstract}

Define the Euler numbers and their normalization by
\[
 \sec t+\tan t=\sum_{n\geq0}E_n\frac{t^n}{n!},\qquad a_n=\frac{E_n}{n!}.
\]
Thus $E_0,E_1,E_2,E_3=1,1,1,2$. Sokal \cite[Conjecture 3.3, p.~14]{Sokal} asks whether the following holds.
\begin{theorem}\label{main}
For every integer $n\geq0$ and $j,k\geq1$,
\[
 (-1)^{n+1}\bigl(a_na_{n+j+k}-a_{n+j}a_{n+k}\bigr)>0.
\]
\end{theorem}
The exponent $n+1$ gives the same sign as the source's integer exponent $n-1$, including $n=0$. The known log-concavity of the even subsequence and log-convexity of the odd subsequence are credited to the source and its references; they are not separate contributions here.

\begin{lemma}\label{errors}
There are real numbers $c_n$ such that
\begin{equation}\label{normal}
 a_n=2(2/\pi)^{n+1}\bigl(1+(-1)^{n+1}c_n\bigr),\qquad
 0<c_n<1,\qquad c_{n+2}\leq c_n/3.
\end{equation}
\end{lemma}
\begin{proof}
Put $c_0=1-\pi/4$. For even $n\geq2$ and odd $n\geq1$, respectively, put
\begin{align*}
 c_n&=\sum_{r\geq0}\left((4r+3)^{-n-1}-(4r+5)^{-n-1}\right),\\
 c_n&=\sum_{r\geq0}(2r+3)^{-n-1}.
\end{align*}
The classical secant and tangent partial fractions, already used in \cite[eqs.~(3.3)--(3.7)]{Sokal}, give \eqref{normal}'s coefficient identity. These series are absolutely convergent for $n\geq1$. At $n=0$ the identity is immediate from $a_0=1$, so no rearrangement of a conditionally convergent series is needed.

For the even expression, telescope a majorant to obtain
\[
 0<c_n\leq3^{-n-1}<1:
 \quad (4r+3)^{-s}-(4r+5)^{-s}
 \leq(4r+3)^{-s}-(4r+7)^{-s},\quad s=n+1.
\]
For $0<B<A\leq1/3$ and integer $s\geq1$, the identity
\[
 (A^2+AB+B^2)(A^s-B^s)-(A^{s+2}-B^{s+2})
 =AB(A+B)(A^{s-1}-B^{s-1})\geq0
\]
and $A^2+AB+B^2\leq1/3$ prove the required contraction term by term.

For odd $n$, comparison with
\[
 (2r+3)^{-2}\leq\frac1{(r+2)(r+3)}
 =\frac1{r+2}-\frac1{r+3}
\]
gives $0<c_n\leq1/2<1$. Each summand contracts by at most $1/9$, hence by at most $1/3$, when $n$ increases by two. Finally, $0<c_0<1$ follows from $0<\pi<4$, and
$c_2\leq1/27<(1-\pi/4)/3=c_0/3$, using the classical bound $\pi<63/20$.
\end{proof}

\begin{lemma}\label{criterion}
If $0<c_n<1$ and $c_{n+2}\leq c_n/3$ for every $n\geq0$, then
$b_n=1+(-1)^{n+1}c_n$ satisfies Theorem~\ref{main}'s inequality with $a$ replaced by $b$.
\end{lemma}
\begin{proof}
Every positive even shift $h$ satisfies $c_{n+h}\leq c_n/3$. Set
$a=c_n$, $b=c_{n+j}$, $c=c_{n+k}$, $d=c_{n+j+k}$ and $X=ad-bc$.
Multiplication of the expanded determinant by $\varepsilon=(-1)^{n+1}$ gives
\[
 L+\varepsilon(-1)^{j+k}X,\qquad
 L=a+(-1)^{j+k}d-(-1)^jb-(-1)^kc.
\]
It suffices to show both $L-X>0$ and $L+X>0$.

If $j,k$ are odd, $L=a+b+c+d$, and
\[
 L-X=a(1-d)+d+b+c+bc>0,\quad
 L+X=b(1-c)+c+a+d+ad>0.
\]
If $j,k$ are even, $b,c\leq a/3$ and $L=a+d-b-c$. Thus
\[
 L-X=(a-b-c)+d(1-a)+bc>0,\qquad
 L+X\geq a/3-a^2/9>0.
\]
If $j$ is odd and $k$ even, $c\leq a/3$, $d\leq b/3$ and $L=a+b-c-d$. For either sign,
\[
 L\pm X\geq\tfrac23(a+b)-\tfrac13ab>0.
\]
The other mixed case follows by exchanging $j,k$.
\end{proof}
\begin{proof}[Proof of Theorem~\ref{main}]
Apply Lemma~\ref{criterion} to Lemma~\ref{errors}. The two products of geometric factors in the determinant are equal and strictly positive, so they preserve its sign.
\end{proof}

\section*{Formal scope and status}
The endpoint \texttt{EulerTuran.source\_conjecture\_338918} has only the source index hypotheses $n\in\mathbb N$ and $j,k\geq1$. The definition \texttt{normalizedEuler} is the real part of the actual complex derivative of $1/\cos z+\tan z$ at zero, divided by $n!$. Analyticity, reality of these coefficients and their genuine power-series representation are proved; no abstract error sequence is substituted for the source object.

The coefficient bridge starts from mathlib's cotangent derivative expansion in the upper half-plane. A dominated limit extends it to real nonintegers: along $a+iy$ the reciprocal-power summands are bounded by $|m+a|^{-n-1}$, a summable majorant for $n\geq1$. The analytic identity $\sec z+\tan z=\cot(\pi/4-z/2)$ then yields the actual coefficients. The zeroth coefficient is treated directly. All infinite sums, parity reindexings, estimates and the finite sign criterion are proved in Lean 4.34.1 with pinned mathlib revision \texttt{d13f23b723b8a846827a245b89c10fc7d3f11612}.

Partial fractions, standard analytic and summation facts, and the source's known subsequence results are credited ingredients. Finite checks and helper lemmas receive no separate credit. No worldwide-priority claim is made.

\begin{thebibliography}{9}
\bibitem{Sokal} A.~D. Sokal, \emph{The Euler and Springer numbers as moment sequences}, Expositiones Mathematicae \textbf{38} (2020), 1--26; arXiv:1804.04498v3 (2018), \href{https://doi.org/10.1016/j.exmath.2018.08.001}{doi:10.1016/j.exmath.2018.08.001}.
\end{thebibliography}
\end{document}
